% CS F407 Artificial Intelligence: final project report template.
% Limit: 2 pages, including everything except references.
% Do not change the font size, margins or spacing.
\documentclass[11pt,a4paper,twocolumn]{article}

\usepackage[margin=2cm,columnsep=0.8cm]{geometry}
\usepackage[T1]{fontenc}
\usepackage{lmodern}
\usepackage{microtype}
\usepackage{amsmath,amssymb}
\usepackage{graphicx}
\usepackage{booktabs}
\usepackage{tabularx}
\usepackage{enumitem}
\usepackage{titlesec}
\usepackage[hidelinks]{hyperref}

\titleformat{\section}{\normalsize\bfseries}{\thesection}{0.6em}{}
\titlespacing*{\section}{0pt}{1.2ex plus 0.3ex}{0.6ex}
\setlist{nosep,leftmargin=*}
\setlength{\parskip}{0.3ex}

\newcommand{\coursename}{CS F407 Artificial Intelligence, AY 2026--27 Semester I}
\newcommand{\teamno}{Team 200 (Team Found)}
\newcommand{\githuburl}{https://github.com/SupercalifragilisticexpialidociousPoppy/turtlebot_platoon}


\begin{document}

\twocolumn[{%
\begin{center}
  {\large\bfseries Robot Platoon: LiDAR-Only Leader Following on TurtleBot3 Robots}\\[1ex]
  \teamno\\
  Arnav Kulshreshtha (2024ADPS0821G), Mohammad Zaid Agha (2024ADPS0840G), Manas Milind Desai (2023A1PS0098G)\\[0.5ex]
  {\small \coursename, BITS Pilani, K.K. Birla Goa Campus}\\[1.5ex]
  \begin{minipage}{0.85\textwidth}
    \small
    \centerline{\bfseries Abstract}\vspace{0.5ex}
    A platoon is a convoy in which each vehicle follows the one ahead at a fixed gap. We
    built a platoon of TurtleBot3 Burger robots in which every follower is an independent
    reactive agent: it uses only its own 2D LiDAR, with no communication, cameras or maps,
    and runs identical code. Each robot isolates the leader in a ±15° window that turns with
    it, rejects spurious returns with a robust clustering step, holds a 30\,cm gap with a PD
    controller, and steers with pure pursuit. In closed-loop simulation against a randomly
    moving target, a gain of $k_p{=}2.0$ kept the worst gap error at 4.1\,cm, inside the
    ±5\,cm requirement, while $k_p{=}1.5$ failed at 5.4\,cm, as predicted by a steady-state
    lag analysis. The same code was deployed on physical robots.
  \end{minipage}
\end{center}
\vspace{1.5ex}
}]

\section{Introduction}
Platooning lets vehicles travel close together under coordinated control, improving
throughput in transport and warehouse logistics. Most systems rely on vehicle-to-vehicle
communication or external tracking. We study the opposite extreme: each robot follows the
object directly ahead using only on-board sensing, so the same program works on any robot in
any position in the line.

\textbf{Input:} a 360-beam laser scan at 5\,Hz (ROS \texttt{LaserScan}).
\textbf{Output:} linear and angular velocity commands (ROS \texttt{Twist}).

\textbf{Contributions:} (i) a LiDAR perception module that tracks a leader through turns and
is robust to single-beam outliers; (ii) a PD gap controller whose gain is derived from the
tolerance specification; (iii) a pure-pursuit steering law with actuator-aware limiting; and
(iv) a Gazebo testbed and a deployment on physical TurtleBot3 robots, with unit tests.

\section{Background and related work}
Each follower is a \emph{simple reflex agent}~\cite{russell2021aima}: it maps the current
percept, plus a small internal state (filtered estimates), directly to an action, without a
world model or planning. For control we use two standard techniques: proportional-derivative
(PD) feedback for the gap, and pure-pursuit path tracking for steering. A central concern in
platoons is \emph{string stability}: disturbances must not grow as they pass down the chain.
The system is built on ROS~2 Humble~\cite{ros2docs} and the ROBOTIS TurtleBot3
platform~\cite{turtlebot3manual}.

\section{Method}
Two ROS~2 nodes run on every follower:
\texttt{scan} $\rightarrow$ \texttt{scan\_filter} $\rightarrow$ \texttt{platoon/target}
$\rightarrow$ \texttt{platoon\_follower} $\rightarrow$ \texttt{cmd\_vel}.
Topic names are relative, so robots can be namespaced.

\textbf{Perception.} Beam angles are wrapped to $(-\pi,\pi]$ (the LDS reports $0$ to $2\pi$,
and without wrapping a $\pm15^\circ$ test keeps only the left half of the cone: 16 of 31
beams). Beams are kept if they lie within $\pm15^\circ$ of a window centre $c$ and have range
in $[0.12, 1.2]$\,m. The centre $c$ follows the last bearing (clamped to $\pm40^\circ$) so
the leader stays in view through turns, and resets to $0$ when the target is lost or closer
than 0.20\,m (to avoid edge-beam cone drift). Let $r_k$ be the $k$-th smallest range
($k{=}3$). The leader's \emph{face} is
$F=\{i : |r_i - r_k| \le 0.10\,\text{m}\}$, accepted if $|F|\ge3$. Then
\begin{equation}
  d=\operatorname{median}_{i\in F} r_i,\qquad
  \beta=\tfrac{1}{|F|}\textstyle\sum_{i\in F}\theta_i .
\end{equation}
Anchoring on $r_k$ rather than the nearest return means a single spurious reading cannot
anchor the cluster. Estimates are smoothed with an exponential moving average
($\alpha{=}0.4$); a detection more than 0.75\,m from the previous target is rejected as a
jump to another object; up to 3 missing frames are bridged before the target is declared
lost; and scan gaps longer than 1.0\,s reset tracker state. The target is published as the
point $(d\cos\beta, d\sin\beta)$.

\textbf{Gap control (PD).} With setpoint $d^\ast{=}0.30$\,m, reverse threshold
$d_{\text{rev}}{=}0.20$\,m, and error $e=d-d^\ast$,
\begin{equation}
  v = \operatorname{clip}\big(k_p e + k_d \dot{\hat d},\; -0.08,\; 0.22\big)\ \text{m/s},
\end{equation}
where $\dot{\hat d}$ is a low-pass-filtered range rate ($\alpha{=}0.3$). On hardware,
distances in $[0.20, 0.30]$\,m form a hold zone ($v{=}0$); above 0.30\,m the follower
drives forward ($v\ge0$, preventing noisy derivative reverse on far targets), and below
0.20\,m it reverses proportionally to $d-0.20$\,m. In the 1D simulation baseline, a symmetric
$\pm1.5$\,cm deadband is used. For a leader moving at constant speed $v_\ell$, a P
controller settles at lag $e_{ss}=v_\ell/k_p$. Requiring $e_{ss}\le5$\,cm at
$v_\ell{=}0.08$\,m/s gives $k_p\ge1.6$; we use $k_p{=}2.0$, $k_d{=}0.5$. We deliberately
omit an integral term: integral windup while a follower is blocked releases as a surge when
it is freed, a disturbance that propagates down the platoon.

\textbf{Steering (pure pursuit).} For target $(x,y)$, the arc through it has curvature
$\kappa=2y/(x^2+y^2)$ and we command $\omega=v\kappa$. At low speed
$\omega=v\kappa\to0$, so $\omega$ is blended towards $k_a\beta$ ($k_a{=}1.4$) as $v$ drops
below 0.08\,m/s, keeping the leader centred while stationary. Forward speed is reduced for
large $|\beta|$, and turning is suppressed below 0.20\,m so close-range backoff runs
straight.

\textbf{Actuation and safety.} (i) If either wheel speed $v\pm\omega b/2$ ($b{=}0.16$\,m)
exceeds 0.22\,m/s, $v$ and $\omega$ are scaled down together, preserving the arc.
(ii) $\omega$ and $v$ are rate limited (2.0\,rad/s$^2$, 0.6\,m/s$^2$): the LiDAR samples its
beams over one 0.2\,s revolution, so fast rotation smears the scan (17° at 1.5\,rad/s) and
corrupts the bearing that produced it. (iii) The robot stops if no target arrives for 0.6\,s
(resetting derivative state) and backs off at $-0.08$\,m/s if closer than 0.16\,m.
(iv) Gains and motor-direction inversion flags are live-tunable via ROS parameters.

\section{Experimental setup}
\textbf{Simulation.} Gazebo~11 with ROS~2 Humble in Docker. A $1.0\times0.5$\,m wall moves
along one axis between 0.6 and 2.5\,m with a seeded random speed profile (up to 0.08\,m/s,
acceleration-limited to 0.05\,m/s$^2$), so every gain setting faces the same disturbance.
Gains were first screened in a kinematic closed-loop replay of this profile with a 5\,Hz
sensor and 2\,mm range noise.

\textbf{Hardware.} TurtleBot3 Burger robots~\cite{turtlebot3manual} (Raspberry~Pi,
OpenCR, LDS-01 LiDAR at 5\,Hz), running identical code on ROS~2 Humble~\cite{ros2docs}. The
leader is driven by keyboard teleoperation. Makeshift height extensions were added to the
TurtleBots to improve target tracking, since all the LiDARs were at the same height.

\textbf{Metrics.} Mean and worst-case absolute gap error after the initial approach (lower is
better), the fraction of time within $\pm5$\,cm (higher is better), and the number of target
losses. \textbf{Baselines:} P-only control ($k_d{=}0$), lower gains, and nearest-return
anchoring.

\section{Results}
Table~\ref{tab:results} shows the gain screening. The pass/fail boundary matches the
$k_p\ge1.6$ prediction: $k_p{=}1.5$ fails by 0.4\,cm, $k_p{=}2.0$ passes. Raising $k_p$ to
3.0 reduces error further, but leaves less stability margin given the 0.2\,s sensor period.
$k_d$ had little effect in this replay because it models no inertia; on hardware it damps
overshoot.

\begin{table}[t]
\centering
\small
\begin{tabular}{lccc}
\toprule
$k_p$, $k_d$ & Mean & Worst & $\pm5$\,cm \\
\midrule
1.0, 0.5 & 3.35 & 7.98 & fail \\
1.5, 0.0 (P only) & 2.56 & 5.38 & fail \\
1.5, 0.5 & 2.45 & 5.39 & fail \\
\textbf{2.0, 0.5 (ours)} & \textbf{2.00} & \textbf{4.12} & \textbf{pass} \\
3.0, 0.5 & 1.58 & 3.58 & pass \\
\bottomrule
\end{tabular}
\caption{Absolute gap error in cm over a 3-minute seeded run, excluding the first 10\,s
(lower is better).}
\label{tab:results}
\end{table}

\textbf{Perception.} With a flat target at 1.0\,m and one spurious 0.08\,m reading injected,
nearest-return anchoring discarded the whole frame, whereas $k$-th-nearest anchoring kept the
lock and reported 1.01\,m from all 31 beams; with no real object, two spikes alone were still
rejected. The perception and control functions are covered by 30 unit tests that run without
ROS.


\section{Discussion and limitations}
\textbf{Sim-to-real gaps.} Several problems only appeared on hardware: the scan-smearing
effect (Gazebo samples all beams at once), a robot whose motors ran in reverse (handled by
\texttt{invert\_linear}), the sensor being about 9\,cm behind the bumper (the setpoint is
measured from the LiDAR), and processes from earlier sessions holding the LiDAR's serial port.

\textbf{Failure modes.} Clustering by range alone cannot separate two objects at the same
distance, such as both posts of a doorway; their mean bearing points at the gap between them.
A window that follows its target can also drift onto a nearby wall during a turn; the jump
threshold and close-range forward lock limit this. Pure pursuit with a 30\,cm lookahead gives
aggressive steering at short range. \textbf{Next steps:} cluster by angular adjacency,
tighten jump rejection, and add feed-forward from an estimate of the leader's speed to
remove the steady-state lag without an integral term.

\section{Conclusion}
A purely reactive, LiDAR-only controller is sufficient for TurtleBot3 platooning: robust
clustering, a PD gap law with a specification-derived gain, and rate-limited pure pursuit hold
the gap within tolerance in simulation, and the same code runs unchanged on every robot.

\section*{Individual member contributions}
{\small
\begin{tabularx}{\linewidth}{@{}l X@{}}
\toprule
Name & Contribution \\
\midrule
Arnav Kulshreshtha & Robot assembly, hardware bring-up and calibration (Milestone~1),
pure-pursuit steering and unified follower node (\texttt{platoon\_follower},
\texttt{control\_math}), dynamic scan tracking, 3-zone stability fixes, launch/test
scripts and unit tests, multi-robot deployment and testing, video; Sections 4 and 6. \\
Mohammad Zaid Agha & Robot assembly and bring-up (Milestone~1), Gazebo testbed
(Milestone~2), kinematic validation and presentation support; Sections 2 and 5. \\
Manas Milind Desai & LiDAR scan filter, PD gap controller and Gazebo moving-wall testbed
(Milestone~2), hardware tuning and safety limits; Sections 1, 3 and 7. \\
\bottomrule
\end{tabularx}}

\smallskip
\noindent All members have read this report and agree with the statement above.

\section*{Generative AI disclosure}
We used Claude (Anthropic; Claude Opus~5 and Opus~5.5) and Google Antigravity (Gemini) to
help write and debug the ROS~2 nodes, unit tests and helper scripts, to troubleshoot the
Docker, Gazebo and robot setup, and to draft and edit this report. We reviewed, tested and
ran all generated code on simulation and hardware, checked all text, and take full
responsibility for both. No LLM is part of the method itself.

\section*{Code availability}
\noindent Code and data:\\
\url{\githuburl}\\
The repository is public and its README contains instructions to reproduce the results.

\begin{thebibliography}{9}
\small
\bibitem{russell2021aima}
S.~Russell and P.~Norvig.
\newblock \emph{Artificial Intelligence: A Modern Approach}, 4th edition.
\newblock Pearson, 2021.

\bibitem{ros2docs}
Open Robotics.
\newblock \emph{ROS~2 Documentation: Humble}.
\newblock \url{https://docs.ros.org/en/humble/}. Accessed November 2026.

\bibitem{turtlebot3manual}
ROBOTIS.
\newblock \emph{TurtleBot3 e-Manual}.
\newblock \url{https://emanual.robotis.com/docs/en/platform/turtlebot3/overview/}.
\newblock Accessed November 2026.
\end{thebibliography}

\end{document}
